\documentclass[11pt]{article}
\usepackage[margin=1in]{geometry}
\usepackage{amsmath, amssymb, amsthm}
\usepackage{enumitem}
\usepackage{hyperref}

\title{Math - Test Sections Questions: Auxiliary Equations with Complex Roots}
\author{Curated Problem Sets}
\date{}

\begin{document}

\maketitle

\tableofcontents
\newpage

% ---------------------------------------------------------
\section{Concept Checklist \& Question Mapping}
% ---------------------------------------------------------
This section outlines all the core concepts, problem types, and tricks extracted from the provided text. Each concept is mapped to its corresponding problems in the Problem Sets.

\begin{itemize}
    \item \textbf{C1: Pure Imaginary Roots (Simple Harmonic)} - Auxiliary equation yields $r = \pm i\beta$. Solution is purely oscillatory. \textit{(Questions: 1, 2, 3)}
    \item \textbf{C2: Complex Roots with Real Part (Damped Harmonic)} - Auxiliary equation yields $r = \alpha \pm i\beta$. Solution involves exponential decay/growth and oscillation. \textit{(Questions: 4, 5, 6, 7)}
    \item \textbf{C3: Initial Value Problems (IVPs) - Pure Imaginary} - Finding specific constants $c_1, c_2$ using $y(0)$ and $y'(0)$ without the Product Rule. \textit{(Questions: 8, 9, 10)}
    \item \textbf{C4: Initial Value Problems (IVPs) - Complex Roots} - Finding specific constants $c_1, c_2$ using $y(0)$ and $y'(0)$. Requires the Product Rule for differentiation. \textit{(Questions: 11, 12, 13, 14, 15)}
    \item \textbf{C5: Higher-Order Equations} - 3rd order or higher ODEs requiring factoring (Rational Root Theorem, synthetic division) into linear and irreducible quadratic factors. \textit{(Questions: 16, 17, 18, 19, 20)}
    \item \textbf{C6: Physical Applications \& Modeling} - Translating physical constants (RLC Circuits) or economic models (Business Cycles) into an ODE and solving the IVP. \textit{(Questions: 21, 22, 23)}
    \item \textbf{C7: Boundary Value Problems (BVPs)} - Conditions specified at two different points (e.g., $t=0$ and $t=L$). Testing for unique, infinite, or no solutions. \textit{(Questions: 24, 25, 26, 27)}
    \item \textbf{C8: Diagnostic Testing ("Find the Flaw")} - Identifying common pitfalls such as missing negative signs, swapping $\alpha$ and $\beta$, or forgetting the Product/Chain rules. \textit{(Questions: 28, 29, 30)}
\end{itemize}

\newpage

% ---------------------------------------------------------
\section{Problem Sets}
% ---------------------------------------------------------

\subsection*{Part A: Pure Imaginary Roots (Simple Harmonic)}
Find the general solution for the following differential equations:
\begin{enumerate}[label=\textbf{\arabic*.}]
    \item $y'' + 25y = 0$
    \item $y'' + 7y = 0$
    \item $4y'' + 36y = 0$
\end{enumerate}

\subsection*{Part B: Complex Roots with Real Part (Damped Harmonic)}
Find the general solution for the following differential equations:
\begin{enumerate}[label=\textbf{\arabic*.}, resume]
    \item $y'' - 4y' + 13y = 0$
    \item $y'' + 6y' + 25y = 0$
    \item $y'' + 2y' + 10y = 0$
    \item $2y'' - 2y' + 5y = 0$
\end{enumerate}

\subsection*{Part C: Initial Value Problems (IVPs) - Pure Imaginary}
Solve the following initial value problems:
\begin{enumerate}[label=\textbf{\arabic*.}, resume]
    \item $y'' + 16y = 0, \quad y(0) = 2, \quad y'(0) = 12$
    \item $y'' + 4y = 0, \quad y(0) = -1, \quad y'(0) = 6$
    \item $9y'' + y = 0, \quad y(0) = 4, \quad y'(0) = -1$
\end{enumerate}

\subsection*{Part D: Initial Value Problems (IVPs) - Complex Roots}
Solve the following initial value problems:
\begin{enumerate}[label=\textbf{\arabic*.}, resume]
    \item $y'' + 4y' + 5y = 0, \quad y(0) = 1, \quad y'(0) = 0$
    \item $y'' - 2y' + 2y = 0, \quad y(0) = 3, \quad y'(0) = 1$
    \item $y'' + 6y' + 13y = 0, \quad y(0) = 0, \quad y'(0) = 8$
    \item $y'' + 2y' + 17y = 0, \quad y(0) = -2, \quad y'(0) = 4$
    \item $y'' - 4y' + 29y = 0, \quad y(0) = 1, \quad y'(0) = 7$
\end{enumerate}

\subsection*{Part E: Higher-Order Equations}
Find the general solution for the following higher-order differential equations:
\begin{enumerate}[label=\textbf{\arabic*.}, resume]
    \item $y''' - y'' + 4y' - 4y = 0$
    \item $y''' + 3y'' + 9y' - 13y = 0$
    \item $y''' + 8y = 0$
    \item $y''' - 2y'' + 2y' = 0$
    \item $y^{(4)} - y = 0$
\end{enumerate}

\subsection*{Part F: Physical Applications \& Modeling}
\begin{enumerate}[label=\textbf{\arabic*.}, resume]
    \item \textbf{RLC Circuit:} An RLC circuit is governed by $LI'' + RI' + \frac{1}{C}I = E'(t)$. Given $L = 1$ H, $R = 4$ $\Omega$, $C = 0.2$ F, and $E(t) = 0$ V. Find the current $I(t)$ if $I(0) = 2$ A and $I'(0) = -4$ A/s.
    \item \textbf{Macroeconomic Business Cycle:} The deviation of a nation's GDP from equilibrium is modeled by $Y'' + 2Y' + 10Y = 0$. Find the trajectory of the GDP, $Y(t)$, after a market shock where $Y(0) = -5$ and $Y'(0) = 0$.
    \item \textbf{RLC Circuit with Constant Voltage:} Given $L = 2$ H, $R = 4$ $\Omega$, $C = \frac{1}{20}$ F, and a constant voltage $E(t) = 50$ V. Find the current $I(t)$ if the initial charge $q(0) = 0$ and initial current $I(0) = 0$. \textit{(Hint: Use $LI'(0) + RI(0) + \frac{q(0)}{C} = E(0)$ to find $I'(0)$).}
\end{enumerate}

\subsection*{Part G: Boundary Value Problems (BVPs)}
Solve the following boundary value problems. State if there is a unique solution, infinitely many solutions, or no solution.
\begin{enumerate}[label=\textbf{\arabic*.}, resume]
    \item $y'' + 9y = 0, \quad y(0) = 0, \quad y(\pi/3) = 0$
    \item $y'' + 25y = 0, \quad y(0) = 1, \quad y(\pi) = 1$
    \item $y'' + 4y = 0, \quad y(0) = 0, \quad y(\pi/4) = 2$
    \item $y'' + y = 0, \quad y(0) = 2, \quad y(\pi/2) = 3$
\end{enumerate}

\subsection*{Part H: Advanced Diagnostic Testing ("Find the Flaw")}
Act as the grader. Identify the specific error, explain why it is an error in one sentence, and provide the correct final solution.
\begin{enumerate}[label=\textbf{\arabic*.}, resume]
    \item \textbf{Task:} Solve $y'' + 36y = 0$. \\
    \textbf{Flawed Solution:} The auxiliary equation is $r^2 + 36 = 0$. Solving for $r$, we get $r^2 = 36 \implies r = \pm 6$. Therefore, the general solution is $y(t) = c_1e^{6t} + c_2e^{-6t}$.
    
    \item \textbf{Task:} Solve $y'' - 8y' + 25y = 0$. \\
    \textbf{Flawed Solution:} The auxiliary equation is $r^2 - 8r + 25 = 0$. Using the quadratic formula: $r = \frac{8 \pm \sqrt{64 - 100}}{2} = \frac{8 \pm \sqrt{-36}}{2} = 4 \pm 3i$. Here, $\alpha = 4$ and $\beta = 3$. The general solution is $y(t) = e^{3t}(c_1 \cos 4t + c_2 \sin 4t)$.
    
    \item \textbf{Task:} Solve the IVP $y'' + 4y' + 13y = 0, \quad y(0) = 2, \quad y'(0) = 0$. \\
    \textbf{Flawed Solution:} $r^2 + 4r + 13 = 0 \implies r = -2 \pm 3i$. General solution: $y(t) = e^{-2t}(c_1 \cos 3t + c_2 \sin 3t)$. Using $y(0) = 2 \implies c_1 = 2$. Now, take the derivative: $y'(t) = e^{-2t}(-3c_1 \sin 3t + 3c_2 \cos 3t)$. Using $y'(0) = 0 \implies e^0(-0 + 3c_2) = 0 \implies 3c_2 = 0 \implies c_2 = 0$. Final answer: $y(t) = 2e^{-2t}\cos 3t$.
\end{enumerate}

\newpage

% ---------------------------------------------------------
\section{Detailed Solutions}
% ---------------------------------------------------------

\subsection*{Part A: Pure Imaginary Roots}
\textbf{1.} Auxiliary equation: $r^2 + 25 = 0 \implies r^2 = -25 \implies r = \pm 5i$. Here $\alpha=0, \beta=5$. \\
\textbf{Solution:} $y(t) = c_1 \cos 5t + c_2 \sin 5t$

\vspace{0.5cm}
\noindent\textbf{2.} Auxiliary equation: $r^2 + 7 = 0 \implies r^2 = -7 \implies r = \pm i\sqrt{7}$. Here $\alpha=0, \beta=\sqrt{7}$. \\
\textbf{Solution:} $y(t) = c_1 \cos(\sqrt{7}t) + c_2 \sin(\sqrt{7}t)$

\vspace{0.5cm}
\noindent\textbf{3.} Divide by 4: $y'' + 9y = 0$. Auxiliary equation: $r^2 + 9 = 0 \implies r = \pm 3i$. \\
\textbf{Solution:} $y(t) = c_1 \cos 3t + c_2 \sin 3t$

\subsection*{Part B: Complex Roots with Real Part}
\textbf{4.} Auxiliary equation: $r^2 - 4r + 13 = 0$. Using quadratic formula: $r = \frac{4 \pm \sqrt{16 - 52}}{2} = \frac{4 \pm \sqrt{-36}}{2} = 2 \pm 3i$. Here $\alpha=2, \beta=3$. \\
\textbf{Solution:} $y(t) = e^{2t}(c_1 \cos 3t + c_2 \sin 3t)$

\vspace{0.5cm}
\noindent\textbf{5.} Auxiliary equation: $r^2 + 6r + 25 = 0$. $r = \frac{-6 \pm \sqrt{36 - 100}}{2} = \frac{-6 \pm \sqrt{-64}}{2} = -3 \pm 4i$. Here $\alpha=-3, \beta=4$. \\
\textbf{Solution:} $y(t) = e^{-3t}(c_1 \cos 4t + c_2 \sin 4t)$

\vspace{0.5cm}
\noindent\textbf{6.} Auxiliary equation: $r^2 + 2r + 10 = 0$. $r = \frac{-2 \pm \sqrt{4 - 40}}{2} = -1 \pm 3i$. Here $\alpha=-1, \beta=3$. \\
\textbf{Solution:} $y(t) = e^{-t}(c_1 \cos 3t + c_2 \sin 3t)$

\vspace{0.5cm}
\noindent\textbf{7.} Auxiliary equation: $2r^2 - 2r + 5 = 0$. $r = \frac{2 \pm \sqrt{4 - 40}}{4} = \frac{2 \pm 6i}{4} = \frac{1}{2} \pm \frac{3}{2}i$. Here $\alpha=1/2, \beta=3/2$. \\
\textbf{Solution:} $y(t) = e^{t/2}\left(c_1 \cos\left(\frac{3}{2}t\right) + c_2 \sin\left(\frac{3}{2}t\right)\right)$

\subsection*{Part C: Initial Value Problems (IVPs) - Pure Imaginary}
\textbf{8.} $r^2 + 16 = 0 \implies r = \pm 4i$. General solution: $y(t) = c_1 \cos 4t + c_2 \sin 4t$. \\
Apply $y(0) = 2 \implies c_1 = 2$. \\
Derivative: $y'(t) = -4c_1 \sin 4t + 4c_2 \cos 4t$. \\
Apply $y'(0) = 12 \implies 4c_2 = 12 \implies c_2 = 3$. \\
\textbf{Solution:} $y(t) = 2\cos 4t + 3\sin 4t$

\vspace{0.5cm}
\noindent\textbf{9.} $r^2 + 4 = 0 \implies r = \pm 2i$. General solution: $y(t) = c_1 \cos 2t + c_2 \sin 2t$. \\
Apply $y(0) = -1 \implies c_1 = -1$. \\
Derivative: $y'(t) = -2c_1 \sin 2t + 2c_2 \cos 2t$. \\
Apply $y'(0) = 6 \implies 2c_2 = 6 \implies c_2 = 3$. \\
\textbf{Solution:} $y(t) = -\cos 2t + 3\sin 2t$

\vspace{0.5cm}
\noindent\textbf{10.} $9r^2 + 1 = 0 \implies r^2 = -1/9 \implies r = \pm \frac{1}{3}i$. General solution: $y(t) = c_1 \cos(t/3) + c_2 \sin(t/3)$. \\
Apply $y(0) = 4 \implies c_1 = 4$. \\
Derivative: $y'(t) = -\frac{1}{3}c_1 \sin(t/3) + \frac{1}{3}c_2 \cos(t/3)$. \\
Apply $y'(0) = -1 \implies \frac{1}{3}c_2 = -1 \implies c_2 = -3$. \\
\textbf{Solution:} $y(t) = 4\cos(t/3) - 3\sin(t/3)$

\subsection*{Part D: Initial Value Problems (IVPs) - Complex Roots}
\textbf{11.} $r^2 + 4r + 5 = 0 \implies r = -2 \pm i$. General solution: $y(t) = e^{-2t}(c_1 \cos t + c_2 \sin t)$. \\
Apply $y(0) = 1 \implies c_1 = 1$. \\
Derivative (Product Rule): $y'(t) = -2e^{-2t}(c_1 \cos t + c_2 \sin t) + e^{-2t}(-c_1 \sin t + c_2 \cos t)$. \\
Apply $y'(0) = 0 \implies -2c_1 + c_2 = 0 \implies -2(1) + c_2 = 0 \implies c_2 = 2$. \\
\textbf{Solution:} $y(t) = e^{-2t}(\cos t + 2\sin t)$

\vspace{0.5cm}
\noindent\textbf{12.} $r^2 - 2r + 2 = 0 \implies r = 1 \pm i$. General solution: $y(t) = e^t(c_1 \cos t + c_2 \sin t)$. \\
Apply $y(0) = 3 \implies c_1 = 3$. \\
Derivative: $y'(t) = e^t(c_1 \cos t + c_2 \sin t) + e^t(-c_1 \sin t + c_2 \cos t)$. \\
Apply $y'(0) = 1 \implies c_1 + c_2 = 1 \implies 3 + c_2 = 1 \implies c_2 = -2$. \\
\textbf{Solution:} $y(t) = e^t(3\cos t - 2\sin t)$

\vspace{0.5cm}
\noindent\textbf{13.} $r^2 + 6r + 13 = 0 \implies r = -3 \pm 2i$. General solution: $y(t) = e^{-3t}(c_1 \cos 2t + c_2 \sin 2t)$. \\
Apply $y(0) = 0 \implies c_1 = 0$. So $y(t) = c_2 e^{-3t} \sin 2t$. \\
Derivative: $y'(t) = c_2(-3e^{-3t}\sin 2t + 2e^{-3t}\cos 2t)$. \\
Apply $y'(0) = 8 \implies 2c_2 = 8 \implies c_2 = 4$. \\
\textbf{Solution:} $y(t) = 4e^{-3t}\sin 2t$

\vspace{0.5cm}
\noindent\textbf{14.} $r^2 + 2r + 17 = 0 \implies r = -1 \pm 4i$. General solution: $y(t) = e^{-t}(c_1 \cos 4t + c_2 \sin 4t)$. \\
Apply $y(0) = -2 \implies c_1 = -2$. \\
Derivative: $y'(t) = -e^{-t}(c_1 \cos 4t + c_2 \sin 4t) + e^{-t}(-4c_1 \sin 4t + 4c_2 \cos 4t)$. \\
Apply $y'(0) = 4 \implies -c_1 + 4c_2 = 4 \implies -(-2) + 4c_2 = 4 \implies 4c_2 = 2 \implies c_2 = 1/2$. \\
\textbf{Solution:} $y(t) = e^{-t}\left(-2\cos 4t + \frac{1}{2}\sin 4t\right)$

\vspace{0.5cm}
\noindent\textbf{15.} $r^2 - 4r + 29 = 0 \implies r = 2 \pm 5i$. General solution: $y(t) = e^{2t}(c_1 \cos 5t + c_2 \sin 5t)$. \\
Apply $y(0) = 1 \implies c_1 = 1$. \\
Derivative: $y'(t) = 2e^{2t}(c_1 \cos 5t + c_2 \sin 5t) + e^{2t}(-5c_1 \sin 5t + 5c_2 \cos 5t)$. \\
Apply $y'(0) = 7 \implies 2c_1 + 5c_2 = 7 \implies 2(1) + 5c_2 = 7 \implies 5c_2 = 5 \implies c_2 = 1$. \\
\textbf{Solution:} $y(t) = e^{2t}(\cos 5t + \sin 5t)$

\subsection*{Part E: Higher-Order Equations}
\textbf{16.} Auxiliary equation: $r^3 - r^2 + 4r - 4 = 0$. Factor by grouping: $r^2(r - 1) + 4(r - 1) = 0 \implies (r^2 + 4)(r - 1) = 0$. Roots are $r = 1$ and $r = \pm 2i$. \\
\textbf{Solution:} $y(t) = c_1e^t + c_2 \cos 2t + c_3 \sin 2t$

\vspace{0.5cm}
\noindent\textbf{17.} Auxiliary equation: $r^3 + 3r^2 + 9r - 13 = 0$. By Rational Root Theorem, test $r=1$: $1 + 3 + 9 - 13 = 0$. So $(r-1)$ is a factor. Polynomial division yields $(r-1)(r^2 + 4r + 13) = 0$. The quadratic roots are $r = -2 \pm 3i$. \\
\textbf{Solution:} $y(t) = c_1e^t + e^{-2t}(c_2 \cos 3t + c_3 \sin 3t)$

\vspace{0.5cm}
\noindent\textbf{18.} Auxiliary equation: $r^3 + 8 = 0$. Sum of cubes factoring: $(r+2)(r^2 - 2r + 4) = 0$. Roots are $r = -2$ and $r = \frac{2 \pm \sqrt{4 - 16}}{2} = 1 \pm i\sqrt{3}$. \\
\textbf{Solution:} $y(t) = c_1e^{-2t} + e^t(c_2 \cos(\sqrt{3}t) + c_3 \sin(\sqrt{3}t))$

\vspace{0.5cm}
\noindent\textbf{19.} Auxiliary equation: $r^3 - 2r^2 + 2r = 0 \implies r(r^2 - 2r + 2) = 0$. Roots are $r = 0$ and $r = 1 \pm i$. \\
\textbf{Solution:} $y(t) = c_1 + e^t(c_2 \cos t + c_3 \sin t)$

\vspace{0.5cm}
\noindent\textbf{20.} Auxiliary equation: $r^4 - 1 = 0 \implies (r^2 - 1)(r^2 + 1) = 0 \implies (r-1)(r+1)(r^2+1) = 0$. Roots are $r = \pm 1$ and $r = \pm i$. \\
\textbf{Solution:} $y(t) = c_1e^t + c_2e^{-t} + c_3 \cos t + c_4 \sin t$

\subsection*{Part F: Physical Applications \& Modeling}
\textbf{21.} Equation: $1I'' + 4I' + \frac{1}{0.2}I = 0 \implies I'' + 4I' + 5I = 0$. \\
Roots: $r^2 + 4r + 5 = 0 \implies r = -2 \pm i$. \\
General solution: $I(t) = e^{-2t}(c_1 \cos t + c_2 \sin t)$. \\
$I(0) = 2 \implies c_1 = 2$. \\
$I'(t) = -2e^{-2t}(c_1 \cos t + c_2 \sin t) + e^{-2t}(-c_1 \sin t + c_2 \cos t)$. \\
$I'(0) = -2c_1 + c_2 = -4 \implies -4 + c_2 = -4 \implies c_2 = 0$. \\
\textbf{Solution:} $I(t) = 2e^{-2t}\cos t$

\vspace{0.5cm}
\noindent\textbf{22.} Equation: $Y'' + 2Y' + 10Y = 0$. \\
Roots: $r^2 + 2r + 10 = 0 \implies r = -1 \pm 3i$. \\
General solution: $Y(t) = e^{-t}(c_1 \cos 3t + c_2 \sin 3t)$. \\
$Y(0) = -5 \implies c_1 = -5$. \\
$Y'(0) = -c_1 + 3c_2 = 0 \implies 5 + 3c_2 = 0 \implies c_2 = -5/3$. \\
\textbf{Solution:} $Y(t) = e^{-t}\left(-5\cos 3t - \frac{5}{3}\sin 3t\right)$

\vspace{0.5cm}
\noindent\textbf{23.} Equation: $2I'' + 4I' + 20I = 0 \implies I'' + 2I' + 10I = 0$. \\
Roots: $r = -1 \pm 3i$. General solution: $I(t) = e^{-t}(c_1 \cos 3t + c_2 \sin 3t)$. \\
$I(0) = 0 \implies c_1 = 0$. So $I(t) = c_2e^{-t}\sin 3t$. \\
Find $I'(0)$: $2I'(0) + 4(0) + 20(0) = 50 \implies 2I'(0) = 50 \implies I'(0) = 25$. \\
$I'(t) = c_2(-e^{-t}\sin 3t + 3e^{-t}\cos 3t)$. \\
$I'(0) = 3c_2 = 25 \implies c_2 = 25/3$. \\
\textbf{Solution:} $I(t) = \frac{25}{3}e^{-t}\sin 3t$

\subsection*{Part G: Boundary Value Problems (BVPs)}
\textbf{24.} General solution: $y(t) = c_1 \cos 3t + c_2 \sin 3t$. \\
$y(0) = 0 \implies c_1 = 0$. So $y(t) = c_2 \sin 3t$. \\
$y(\pi/3) = c_2 \sin(\pi) = 0 \implies c_2(0) = 0$. This is true for any $c_2$. \\
\textbf{Solution:} Infinitely many solutions of the form $y(t) = c_2 \sin 3t$.

\vspace{0.5cm}
\noindent\textbf{25.} General solution: $y(t) = c_1 \cos 5t + c_2 \sin 5t$. \\
$y(0) = 1 \implies c_1 = 1$. So $y(t) = \cos 5t + c_2 \sin 5t$. \\
$y(\pi) = \cos(5\pi) + c_2 \sin(5\pi) = -1 + 0 = -1$. \\
But the boundary condition requires $y(\pi) = 1$. Contradiction ($-1 \neq 1$). \\
\textbf{Solution:} No solution.

\vspace{0.5cm}
\noindent\textbf{26.} General solution: $y(t) = c_1 \cos 2t + c_2 \sin 2t$. \\
$y(0) = 0 \implies c_1 = 0$. So $y(t) = c_2 \sin 2t$. \\
$y(\pi/4) = c_2 \sin(\pi/2) = c_2(1) = 2 \implies c_2 = 2$. \\
\textbf{Solution:} Unique solution $y(t) = 2\sin 2t$.

\vspace{0.5cm}
\noindent\textbf{27.} General solution: $y(t) = c_1 \cos t + c_2 \sin t$. \\
$y(0) = 2 \implies c_1 = 2$. So $y(t) = 2\cos t + c_2 \sin t$. \\
$y(\pi/2) = 2\cos(\pi/2) + c_2 \sin(\pi/2) = 0 + c_2(1) = 3 \implies c_2 = 3$. \\
\textbf{Solution:} Unique solution $y(t) = 2\cos t + 3\sin t$.

\subsection*{Part H: Advanced Diagnostic Testing ("Find the Flaw")}
\textbf{28. Flaw:} The student missed the negative sign when moving 36 to the other side; $r^2 = -36$, so the roots are imaginary ($\pm 6i$), not real. \\
\textbf{Correct Solution:} $y(t) = c_1 \cos 6t + c_2 \sin 6t$.

\vspace{0.5cm}
\noindent\textbf{29. Flaw:} The student swapped the real part $\alpha$ and the imaginary part $\beta$ in the general solution formula. \\
\textbf{Correct Solution:} $y(t) = e^{4t}(c_1 \cos 3t + c_2 \sin 3t)$.

\vspace{0.5cm}
\noindent\textbf{30. Flaw:} The student forgot to use the Product Rule when differentiating $y(t) = e^{-2t}(c_1 \cos 3t + c_2 \sin 3t)$. \\
\textbf{Correct Solution:} $y'(t) = -2e^{-2t}(c_1 \cos 3t + c_2 \sin 3t) + e^{-2t}(-3c_1 \sin 3t + 3c_2 \cos 3t)$. Setting $y'(0) = 0$ gives $-2c_1 + 3c_2 = 0 \implies -4 + 3c_2 = 0 \implies c_2 = 4/3$. Final answer: $y(t) = e^{-2t}\left(2\cos 3t + \frac{4}{3}\sin 3t\right)$.

\newpage

% ---------------------------------------------------------
\section{Concept Checklist \& Question Mapping (Appendix)}
% ---------------------------------------------------------
For quick reference and review, here is the mapping of concepts to their respective problem numbers:

\begin{itemize}
    \item \textbf{C1: Pure Imaginary Roots (Simple Harmonic)} $\rightarrow$ Questions 1, 2, 3
    \item \textbf{C2: Complex Roots with Real Part (Damped Harmonic)} $\rightarrow$ Questions 4, 5, 6, 7
    \item \textbf{C3: Initial Value Problems (IVPs) - Pure Imaginary} $\rightarrow$ Questions 8, 9, 10
    \item \textbf{C4: Initial Value Problems (IVPs) - Complex Roots} $\rightarrow$ Questions 11, 12, 13, 14, 15
    \item \textbf{C5: Higher-Order Equations} $\rightarrow$ Questions 16, 17, 18, 19, 20
    \item \textbf{C6: Physical Applications \& Modeling} $\rightarrow$ Questions 21, 22, 23
    \item \textbf{C7: Boundary Value Problems (BVPs)} $\rightarrow$ Questions 24, 25, 26, 27
    \item \textbf{C8: Diagnostic Testing ("Find the Flaw")} $\rightarrow$ Questions 28, 29, 30
\end{itemize}

\end{document}